\documentclass{ximera}

\begin{document}

\title{Homework 5 -- Householder QR and the SVD (Python)}

\begin{abstract}
Implement Householder QR factorization and construct the singular
value decomposition (SVD) using the eigendecomposition of $A^TA$.

Due Date: Wednesday, November 11th, 11:59 pm
\end{abstract}

\maketitle

In this assignment, you will implement (from scratch, with guidance)
an important QR factorization algorithm and use it to explore the
relationship between QR factorizations, eigenvalue decompositions,
and the singular value decomposition. You will complete two Python
functions. Graduate students will complete an additional third
function.

Throughout the assignment, assume
\[
A\in\mathbb{R}^{m\times n}, \qquad m\geq n.
\]

\begin{center}
\begin{tabular}{|c|l|l|l|}
\hline
Part & Function & Algorithm & Output shapes\\
\hline
1 & \texttt{householder\_qr} &
Householder reflections (full QR) &
$Q:m\times m,\ R:m\times n$\\
\hline
2 & \texttt{svd\_reduced} &
Reduced SVD using $A^TA$ &
$U:m\times n,\ S:n\times n,\ V:n\times n$\\
\hline
3 & \texttt{svd\_full} &
Full SVD using $A^TA$ &
$U:m\times m,\ S:m\times n,\ V:n\times n$\\
\hline
\end{tabular}
\end{center}

The exercises are designed to build on one another. In particular,
the full QR factorization will be used to understand orthogonal
transformations, while the SVD implementations will use Python's
built-in NumPy QR function to complete orthonormal bases.

\section*{Programming requirements and restrictions}

Each function must be saved in its own Python file, with the exact
filename specified in the table above.

You should use NumPy for matrix operations. At the beginning of
each implementation file, include:

\begin{verbatim}
import numpy as np
\end{verbatim}

\paragraph{Allowed operations.}

Inside \texttt{householder\_qr.py}, you may use elementary NumPy
operations, including:

\begin{itemize}
    \item Matrix indexing and slicing.
    \item Basic arithmetic operations: \texttt{+}, \texttt{-},
    \texttt{*}, \texttt{/}, and matrix multiplication \texttt{@}.
    \item Transpose using \texttt{.T}.
    \item \texttt{np.zeros}, \texttt{np.eye}, and \texttt{np.shape}.
    \item \texttt{np.linalg.norm} (applied to vectors).
    \item \texttt{np.sqrt}, \texttt{np.maximum}, \texttt{np.abs},
    and \texttt{np.sign}.
    \item Basic control structures (\texttt{for}, \texttt{if},
    and \texttt{while}).
\end{itemize}

Inside the SVD implementations, you may additionally use:

\begin{itemize}
    \item NumPy's built-in \texttt{np.linalg.eigh} function to
    calculate the eigenvalues and eigenvectors of $A^TA$.
    \item NumPy's built-in \texttt{np.linalg.qr} function to
    complete an orthonormal basis for the left singular vectors.
    \item \texttt{np.argsort} to sort eigenvalues.
    \item \texttt{np.finfo(float).eps} to obtain machine precision.
\end{itemize}

The \texttt{np.linalg.qr} function is being used in the SVD
implementations only to complete an orthonormal basis. You must
still implement the SVD construction yourself.

When completing an orthonormal basis, use the complete QR
factorization:

\begin{verbatim}
Q, R = np.linalg.qr(A, mode="complete")
\end{verbatim}

This ensures that $Q$ contains all the required orthonormal columns.

\paragraph{Forbidden operations.}

You may NOT use the following functions inside your implementation
functions:

\begin{itemize}
    \item \texttt{np.linalg.svd}, \texttt{np.linalg.eigvals},
    \texttt{np.linalg.eigvalsh}, \texttt{scipy.linalg.svd}.
    \item \texttt{scipy.linalg.orth}, \texttt{scipy.linalg.null\_space}.
    \item \texttt{np.linalg.cholesky}, \texttt{np.linalg.inv},
    and \texttt{np.linalg.pinv}.
    \item \texttt{np.linalg.solve} and other routines that directly
    solve linear systems.
\end{itemize}

The built-in \texttt{np.linalg.eigh} function is allowed in both
SVD implementations to compute the eigendecomposition of $A^TA$.

The built-in \texttt{np.linalg.qr} function is allowed in both
SVD implementations to complete the orthonormal basis for $U$.

However, you may NOT use NumPy's built-in \texttt{np.linalg.qr}
function inside \texttt{householder\_qr.py}.

\section*{Provided files}

The starter folder contains the following files.

\begin{center}
\begin{tabular}{|p{4cm}|p{9cm}|}
\hline
File(s) & Purpose\\
\hline
\texttt{householder\_qr.py} &
Student implementation of the full QR factorization using
Householder reflections.\\
\hline
\texttt{svd\_reduced.py} &
Student implementation of the reduced SVD using the eigendecomposition
of $A^TA$.\\
\hline
\texttt{svd\_full.py} &
Student implementation of the full SVD using the eigendecomposition
of $A^TA$.\\
\hline
\texttt{check\_qr.py} &
Prints diagnostic information for a QR factorization.\\
\hline
\texttt{check\_svd.py} &
Prints diagnostic information for an SVD.\\
\hline
\texttt{run\_tests.py} &
Runs a collection of tests on the required implementations and
prints PASS or FAIL for each test.\\
\hline
\end{tabular}
\end{center}

\section*{Testing your implementations}

You are provided with two diagnostic functions and a test script.
You should use them throughout the assignment, not just at the end.

\paragraph{QR diagnostics.}

The function \texttt{check\_qr(A,Q,R)} checks:

\begin{enumerate}
    \item The dimensions of $Q$ and $R$.
    \item The relative reconstruction error
    \[
    \frac{\|A-QR\|}{\|A\|}.
    \]
    \item The orthogonality of $Q$:
    \[
    \|Q^TQ-I\|.
    \]
    \item The triangular structure of $R$.
\end{enumerate}

\paragraph{SVD diagnostics.}

The function \texttt{check\_svd(A,U,S,V)} checks:

\begin{enumerate}
    \item The dimensions of $U$, $S$, and $V$.
    \item The relative reconstruction error
    \[
    \frac{\|A-USV^T\|}{\|A\|}.
    \]
    \item The orthogonality of $U$ and $V$.
    \item The agreement between your computed singular values and
    NumPy's built-in \texttt{np.linalg.svd}.
\end{enumerate}

Note that your SVD functions return $V$, not $V^T$.

\paragraph{Running the tests.}

Once your functions are saved in the same folder as the provided
scripts, open a terminal or Python environment in that folder
and run:

\begin{verbatim}
python run_tests.py
\end{verbatim}

The test script checks your implementations on several matrices,
including:

\begin{itemize}
    \item Random tall matrices.
    \item Random square matrices.
    \item Rank-deficient matrices (for the SVD implementations).
\end{itemize}

Your goal is to obtain PASS for every test associated with your
implementation.

If a test fails, use the diagnostic functions to identify which
property is not being satisfied. A small reconstruction error alone
does not guarantee a correct factorization; orthogonality and
matrix dimensions must also be correct.

You should also consider testing your SVD implementations on the
zero matrix and nearly rank-deficient matrices. These are useful
additional tests of your numerical-rank calculations and special
case handling.

\section{Householder QR}

\paragraph{Background.}

Householder QR is a method for computing a QR factorization.
Instead of orthogonalizing individual columns, it uses orthogonal
reflections to eliminate the entries below the diagonal.

A Householder reflector has the form
\[
H=I-2vv^T,
\]
where $\|v\|_2=1$.

At step $k$, consider the vector
\[
x=R(k:m,k).
\]

Construct
\[
v=\operatorname{sign}(x_1)\|x\|_2e_1+x,
\]
and normalize it:
\[
v\leftarrow\frac{v}{\|v\|_2}.
\]

Use the convention $\operatorname{sign}(0)=1$.

In Python, you should explicitly implement this convention when
calculating the sign.

The corresponding reflection maps $x$ to a vector whose entries
below the first position are zero. This eliminates the entries
below the diagonal in column $k$.

Rather than forming the entire matrix $H$, apply the reflection
directly to the trailing block:
\[
R(k:m,k:n)
\leftarrow
R(k:m,k:n)-2v\bigl(v^TR(k:m,k:n)\bigr).
\]

In Python, matrix multiplication should use the \texttt{@}
operator, and the transpose of a matrix or vector should use
\texttt{.T}.

\paragraph{Task.}

Complete the provided function \texttt{householder\_qr.py}.

Your function should return the full QR factorization:
\[
Q\in\mathbb{R}^{m\times m},
\qquad
R\in\mathbb{R}^{m\times n}.
\]

Follow the three-step structure in the starter file.

\begin{enumerate}
    \item Initialize $R=A$ and $Q=I_m$. Create a matrix to
    store the Householder vectors.

    \item Apply the Householder reflections to $R$.

    For each column $k=0,\ldots,n-1$:

    \begin{itemize}
        \item Extract the portion of the column that must be
        eliminated:
        \[
        x=R(k:m,k).
        \]

        \item Skip the reflection if the extracted vector is
        already zero.

        \item Construct and normalize the Householder vector.

        \item Store the vector in the appropriate portion of
        the storage matrix.

        \item Apply the reflection directly to the trailing
        block of $R$.
    \end{itemize}

    \item Construct $Q$ by backward accumulation: starting with
    $Q=I_m$, apply the stored Householder reflections in reverse
    order.

    \item Set the entries strictly below the diagonal of $R$
    to exactly zero before returning.
\end{enumerate}

Remember that Python uses zero-based indexing, so the first
row and column have index 0.

Do not explicitly construct the individual $m\times m$
Householder reflection matrices.

Your final matrices must satisfy
\[
A=QR,\qquad Q^TQ=QQ^T=I_m.
\]

\begin{exercise}

\begin{enumerate}

    \item Why is the sign in
    \[
    v=\operatorname{sign}(x_1)\|x\|_2e_1+x
    \]
    chosen this way?

    \item Run \texttt{check\_qr} on a random $50\times10$ matrix.
    Report the reconstruction error and orthogonality error.

    \item Compare the results of your Householder QR
    implementation with NumPy's built-in \texttt{np.linalg.qr}
    function. Are the factors identical? Explain.

\end{enumerate}

\end{exercise}

\section{Reduced Singular Value Decomposition}

\paragraph{Background.}

The singular value decomposition (SVD) of a matrix is
\[
A=USV^T.
\]

For the reduced SVD, the matrices have dimensions
\[
U\in\mathbb{R}^{m\times n},\qquad
S\in\mathbb{R}^{n\times n},\qquad
V\in\mathbb{R}^{n\times n}.
\]

The columns of $U$ and $V$ are orthonormal, and $S$ is diagonal
with nonnegative singular values in descending order:
\[
\sigma_1\geq\sigma_2\geq\cdots\geq\sigma_n\geq0.
\]

We will construct the SVD using the eigenvalue decomposition
of $A^TA$:
\[
A^TA=V\Lambda V^T.
\]

The singular values are related to the eigenvalues by
\[
\sigma_i=\sqrt{\lambda_i}.
\]

For nonzero singular values, the corresponding left singular
vectors are calculated using
\[
u_i=\frac{Av_i}{\sigma_i}.
\]

\paragraph{Completing the left singular vectors.}

The formula above cannot be used when $\sigma_i=0$.

Furthermore, when $A$ is rank deficient, the left singular
vectors corresponding to zero singular values must still be
orthonormal and orthogonal to the previously calculated columns.

We will use NumPy's built-in \texttt{np.linalg.qr} function to
complete these vectors to an orthonormal basis.

This use of \texttt{qr} does not calculate the SVD. It is used
only to complete the left singular vectors after the nonzero
singular vectors have been calculated.

Use the \texttt{mode="complete"} option so that the QR
factorization returns a full orthogonal matrix.

\paragraph{Task.}

Complete the provided function \texttt{svd\_reduced.py}.

Follow the suggested outline in the starter file.

In particular:

\begin{enumerate}

    \item Use NumPy's built-in \texttt{np.linalg.eigh} function to
    calculate the eigenvalues and eigenvectors of $A^TA$:
    \[
    A^TA=V\Lambda V^T.
    \]

    \item Sort the eigenvalues in descending order and rearrange
    the columns of $V$ accordingly.

    \item Calculate the singular values using
    \[
    \sigma_i=\sqrt{\max(\lambda_i,0)}.
    \]

    \item Determine the numerical rank using
    \[
    \text{tol}=
    \sqrt{\max(m,n)\varepsilon_{\mathrm{mach}}}\,\sigma_1.
    \]

    The numerical rank is
    \[
    r=\#\{i:\sigma_i>\text{tol}\}.
    \]

    Set the remaining singular values to exactly zero.

    In Python, machine precision can be obtained using
    \texttt{np.finfo(float).eps}.

    \item Calculate the first $r$ left singular vectors using
    \[
    u_i=\frac{Av_i}{\sigma_i}.
    \]

    \item Complete the remaining columns of $U$ using NumPy's
    built-in \texttt{np.linalg.qr} function.

    Specifically, apply a complete QR factorization to the
    matrix containing the first $r$ calculated left singular
    vectors. Retain the first $n$ columns of the resulting
    orthogonal matrix.

    \item Handle the special case $r=0$ by taking the first
    $n$ columns of $I_m$.

    \item Return
    \[
    S=\operatorname{diag}(\sigma_1,\ldots,\sigma_n).
    \]

\end{enumerate}

\paragraph{Important numerical consideration.}

The matrix $A^TA$ has condition number approximately equal to
the square of the condition number of $A$:
\[
\operatorname{cond}(A^TA)=\operatorname{cond}(A)^2.
\]

Consequently, very small singular values may be lost when
computing the eigenvalues of $A^TA$.

This is why we use a numerical rank tolerance rather than
assuming that every computed eigenvalue is positive and meaningful.

\begin{exercise}

\begin{enumerate}

    \item Why do we sort the eigenvalues in descending order?

    \item Why can we not directly calculate
    $u_i=Av_i/\sigma_i$ when $\sigma_i=0$?

    \item Explain why the eigenvalues of $A^TA$ are nonnegative
    in exact arithmetic.

    \item Run \texttt{check\_svd} on a random $8\times5$ matrix
    and a rank-deficient matrix. Report your results.

\end{enumerate}

\end{exercise}

\section{GRAD STUDENTS ONLY: Full Singular Value Decomposition}

\paragraph{Background.}

The full SVD has the same factorization
\[
A=USV^T,
\]
but now
\[
U\in\mathbb{R}^{m\times m},\qquad
S\in\mathbb{R}^{m\times n},\qquad
V\in\mathbb{R}^{n\times n}.
\]

Thus, $U$ is a square orthogonal matrix, and $S$ is a rectangular
diagonal matrix.

The full SVD contains the same singular values as the reduced SVD.
The difference is that the full SVD includes all $m$ columns
of $U$.

\paragraph{Task.}

Complete the provided function \texttt{svd\_full.py}.

This implementation follows essentially the same procedure as
the reduced SVD. The main difference is that the full SVD retains
all $m$ columns of $U$.

Follow the suggested outline in the starter file.

In particular:

\begin{enumerate}

    \item Calculate the eigenvalues and eigenvectors of $A^TA$,
    sort them in descending order, and rearrange the eigenvectors
    accordingly.

    \item Calculate the singular values and determine the
    numerical rank using the tolerance specified in Part 2.

    \item Calculate the first $r$ columns of $U$ using
    \[
    u_i=\frac{Av_i}{\sigma_i}.
    \]

    \item Complete $U$ using NumPy's built-in \texttt{np.linalg.qr}
    function. Apply a complete QR factorization to the matrix
    containing the first $r$ left singular vectors and retain
    all $m$ columns of the resulting orthogonal matrix.

    \item Handle the special case $r=0$ by setting
    \[
    U=I_m.
    \]

    \item Construct the rectangular matrix $S$ by embedding
    the diagonal matrix of singular values into an $m\times n$
    zero matrix:
    \[
    S=
    \begin{bmatrix}
    \operatorname{diag}(\sigma_1,\ldots,\sigma_n)\\
    0
    \end{bmatrix}.
    \]

\end{enumerate}

Your final matrices must satisfy
\[
A=USV^T,\qquad U^TU=I_m,\qquad V^TV=I_n.
\]

\begin{exercise}

\begin{enumerate}

    \item What is the difference between the reduced and full SVD?

    \item Explain why the additional columns of $U$ are
    orthogonal to the range of $A$.

    \item What is the relationship between the additional
    columns of $U$ and the null space of $A^T$?

    \item Run \texttt{check\_svd} on a random $8\times5$ matrix
    and a zero matrix. Verify that your results satisfy the
    defining properties of the full SVD.

\end{enumerate}

\end{exercise}

\section*{Final testing and submission}

Before submitting your assignment, make sure that all three
implementation functions are saved in the correct folder.

Run the provided script:

\begin{verbatim}
python run_tests.py
\end{verbatim}

Your implementation should pass all the tests associated with
the required functions.

The provided test script checks:

\begin{itemize}
    \item Random tall matrices.
    \item Random square matrices.
    \item Rank-deficient matrices for the SVD implementations.
\end{itemize}

You should also test your SVD implementations on the zero matrix
and nearly rank-deficient matrices. These are particularly
important for checking the numerical-rank calculation and
special case handling.

If a test fails, use \texttt{check\_qr} or \texttt{check\_svd} to help
identify the source of the error.

\paragraph{Files to submit.}

Submit the following three Python function files:

\begin{enumerate}
    \item \texttt{householder\_qr.py}
    \item \texttt{svd\_reduced.py}
    \item \texttt{svd\_full.py}
\end{enumerate}

Also include a text file or PDF containing the answers to the
conceptual questions in Parts 1--3.

You do not need to submit the provided checking scripts.

\paragraph{Important.}

Your code will be evaluated using the provided tests and
additional test matrices that may not appear in the starter
files.

Passing the provided tests is necessary, but your implementation
should work for general matrices satisfying the assumptions
specified in each exercise.

\end{document}
